\begin{filecontents*}{refs.bib}
@article{li2025questbench,
  title={QuestBench: Can LLMs ask the right question to acquire information in reasoning tasks?},
  author={Li, Belinda Z. and Kim, Been and Wang, Zi},
  journal={arXiv preprint arXiv:2503.22674},
  year={2025}
}

@inproceedings{toles2025factualclarification,
  title={Learning and Evaluating Factual Clarification Question Generation Without Examples},
  author={Toles, Matthew and Huang, Yukun and Yu, Zhou},
  booktitle={Proceedings of the Fourth Workshop on Generation, Evaluation and Metrics},
  year={2025},
  pages={200--211}
}

@article{kuhn2023clam,
  title={CLAM: Selective Clarification for Ambiguous Questions with Generative Language Models},
  author={Kuhn, Lorenz and Gal, Yarin and Farquhar, Sebastian},
  journal={arXiv preprint arXiv:2212.07769},
  year={2023}
}

@inproceedings{lee-etal-2023-asking,
  title = {Asking Clarification Questions to Handle Ambiguity in Open-Domain QA},
  author = {Lee, Dongryeol and Kim, Segwang and Lee, Minwoo and Lee, Hwanhee and Park, Joonsuk and Lee, Sang-Woo and Jung, Kyomin},
  booktitle = {Findings of the Association for Computational Linguistics: EMNLP 2023},
  year = {2023},
  pages = {11526--11544},
  address = {Singapore},
  publisher = {Association for Computational Linguistics},
  doi = {10.18653/v1/2023.findings-emnlp.772},
  url = {https://aclanthology.org/2023.findings-emnlp.772/}
}

@inproceedings{zhang-etal-2024-clamber,
  title = {CLAMBER: A Benchmark of Identifying and Clarifying Ambiguous Information Needs in Large Language Models},
  author = {Zhang, Tong and Qin, Peixin and Deng, Yang and Huang, Chen and Lei, Wenqiang and Liu, Junhong and Jin, Dingnan and Liang, Hongru and Chua, Tat-Seng},
  booktitle = {Proceedings of the 62nd Annual Meeting of the Association for Computational Linguistics},
  year = {2024},
  pages = {10746--10766},
  publisher = {Association for Computational Linguistics},
  doi = {10.18653/v1/2024.acl-long.578},
  url = {https://aclanthology.org/2024.acl-long.578/}
}

@article{chi2024clarinet,
  title={CLARINET: Augmenting Language Models to Ask Clarification Questions for Retrieval},
  author={Chi, Yizhou and Lin, Jessy and Lin, Kevin and Klein, Dan},
  journal={arXiv preprint arXiv:2405.15784},
  year={2024}
}

@article{wu2025clarifycoder,
  title={Can Code Language Models Learn Clarification-Seeking Behaviors?},
  author={Wu, Jie JW and Chaudhary, Manav and Abrahamyan, Davit and Khaku, Arhaan and Wei, Anjiang and Fard, Fatemeh H.},
  journal={arXiv preprint arXiv:2504.16331},
  year={2025}
}

@article{andukuri2024stargate,
  title={STaR-GATE: Teaching Language Models to Ask Clarifying Questions},
  author={Andukuri, Chinmaya and Fr{\"a}nken, Jan-Philipp and Gerstenberg, Tobias and Goodman, Noah D.},
  journal={arXiv preprint arXiv:2403.19154},
  year={2024}
}

@article{zhang2025modeling,
  title={Modeling Future Conversation Turns to Teach LLMs to Ask Clarifying Questions},
  author={Zhang, Michael J.Q. and Knox, W. Bradley and Choi, Eunsol},
  journal={arXiv preprint arXiv:2410.13788},
  year={2025}
}

@inproceedings{tafjord2021proofwriter,
  title={ProofWriter: Generating Implications, Proofs, and Abductive Statements over Natural Language},
  author={Tafjord, Oyvind and Dalvi, Bhavana and Clark, Peter},
  booktitle={Findings of the Association for Computational Linguistics: EMNLP 2021},
  year={2021}
}

@techreport{mcdermott1998pddl,
  title={PDDL - The Planning Domain Definition Language},
  author={McDermott, Drew and Ghallab, Malik and Howe, Adele and Knoblock, Craig and Ram, Ashwin and Veloso, Manuela and Weld, Daniel and Wilkins, David},
  institution={AIPS-98 Planning Committee},
  year={1998}
}

@article{helmert2006fastdownward,
  title={The Fast Downward Planning System},
  author={Helmert, Malte},
  journal={Journal of Artificial Intelligence Research},
  volume={26},
  pages={191--246},
  year={2006}
}

@inproceedings{hu-etal-2023-wont,
  title = {Won't Get Fooled Again: Answering Questions with False Premises},
  author = {Hu, Shengding and Luo, Yifan and Wang, Huadong and Cheng, Xingyi and Liu, Zhiyuan and Sun, Maosong},
  booktitle = {Proceedings of the 61st Annual Meeting of the Association for Computational Linguistics (Volume 1: Long Papers)},
  year = {2023},
  url = {https://aclanthology.org/2023.acl-long.309/},
  doi = {10.18653/v1/2023.acl-long.309},
  pages = {5626--5643}
}
\end{filecontents*}

\documentclass{article}

\IfFileExists{neurips_2025.sty}{
  \usepackage[preprint]{neurips_2025}
}{
  \usepackage[margin=1in]{geometry}
  \usepackage[numbers]{natbib}
  \usepackage{times}
  \usepackage{microtype}
}

\usepackage[T1]{fontenc}
\usepackage[utf8]{inputenc}
\usepackage{url}
\usepackage{booktabs}
\usepackage{amsmath}
\usepackage{amssymb}
\usepackage{graphicx}
\usepackage{xcolor}
\usepackage{multirow}
\usepackage{array}
\usepackage{hyperref}
\usepackage{enumitem}

\hypersetup{colorlinks=true,citecolor=blue,linkcolor=blue,urlcolor=blue}

\title{When Does Clarification Supervision Transfer to Formal Reasoning?\\A Controlled Pilot of Validator-Clean vs. Matched Noisy Missing-Fact Tuning}

\author{Anonymous Author(s)}

\begin{document}

\maketitle

\begin{abstract}
Clarification-seeking behavior has been studied in open-domain QA, retrieval, dialogue, and code, but those settings only weakly isolate whether transfer fails because models cannot recognize missing information or because the supervision itself is noisy. We revisit that question in mechanically checkable reasoning domains by comparing three token-matched fine-tuning conditions for Qwen2.5-3B-Instruct with QLoRA: answer-only supervision, matched noisy clarification supervision, and validator-clean clarification supervision. The original study design targeted both ProofWriter logic and Blocksworld planning, but the executed artifacts support a narrower interpretation: the clearest signal is in planning. On the pre-registered seed, validator-clean tuning reaches 0.636 question accuracy on QuestBench Planning-Q, compared with 0.614 for answer-only and 0.016 for noisy supervision, while also achieving 0.450 question accuracy on a held-out synthetic clean planning set. However, logic-side clarification accuracy remains 0.000 in the main run, the pre-registered human audit was not executed, and a longer-horizon residual-noise check reports a planning instability rate of 1.000. We therefore present the result as a planning-dominant pilot rather than confirmatory evidence: clean supervision can help, but the current artifacts do not yet establish robust cross-domain transfer to formal reasoning.
\end{abstract}

\section{Introduction}
Large language models can be taught to ask clarifying questions, yet the literature usually studies settings where the missing information and the value of a clarification are only partially grounded \citep{kuhn2023clam,lee-etal-2023-asking,zhang-etal-2024-clamber,chi2024clarinet,andukuri2024stargate,wu2025clarifycoder}. Formal reasoning offers a cleaner target. In theorem proving and planning, a single hidden fact can change entailment or solvability, and that change can be checked mechanically.

This paper asks a narrow question: under a fixed small-model fine-tuning budget, does clarification tuning transfer better when the synthetic supervision is validator-clean rather than merely plausible? We build directly on the proposal in this workspace, but the final paper reflects the executed artifacts rather than the intended two-domain study. The implementation still constructs both logic and planning data, yet the strongest evidence is planning-side: logic question accuracy is zero in the main run, the human audit was not completed, and one residual-noise check reveals strong horizon sensitivity for planning.

Our contributions are:
\begin{itemize}[leftmargin=*]
  \item We operationalize matched clean and noisy missing-fact supervision for ProofWriter-style logic and Blocksworld planning, using executable validators instead of subjective labeling.
  \item We run a token-matched three-condition comparison on Qwen2.5-3B-Instruct with QLoRA, keeping the total training budget near 25.5k tokens per condition.
  \item We show a planning-side pilot result: on the main seed, validator-clean supervision improves QuestBench Planning-Q question accuracy to 0.636 versus 0.614 for answer-only and 0.016 for noisy supervision, but the broader formal-reasoning claim remains unconfirmed because logic transfer is flat and the pre-registered human audit is missing.
\end{itemize}

Section~\ref{sec:related} reviews prior clarification and formal reasoning work. Section~\ref{sec:method} describes the executed data construction and training setup. Section~\ref{sec:experiments} reports the planning-dominant empirical results and exploratory multi-seed summaries. Section~\ref{sec:discussion} discusses what can and cannot be concluded from the current artifacts.

\section{Related Work}
\label{sec:related}
Clarification has been studied broadly in ambiguous QA and information-seeking benchmarks. CLAM \citep{kuhn2023clam}, CAmbigNQ \citep{lee-etal-2023-asking}, and CLAMBER \citep{zhang-etal-2024-clamber} study when models should ask follow-up questions in open-domain or ambiguous information settings. These papers establish the importance of clarification, but they do not isolate the effect of supervision cleanliness with mechanically checked missing facts.

Several works study how to induce clarification-seeking behavior from synthetic or structured supervision. CLARINET \citep{chi2024clarinet} targets retrieval, STaR-GATE \citep{andukuri2024stargate} teaches question-asking behavior by iterative self-improvement, and recent work extends clarification supervision to dialogue and code \citep{zhang2025modeling,wu2025clarifycoder}. The closest prior work to our setup is \citet{toles2025factualclarification}, which creates automatically evaluable factual clarification data without human examples. Our setting differs in two ways: we focus on formal reasoning rather than factual QA, and we compare validator-clean supervision directly against matched noisy supervision.

Our benchmarks rely on formal substrates rather than new task definitions. ProofWriter supplies natural-language reasoning instances with proof traces \citep{tafjord2021proofwriter}. Planning is grounded in Blocksworld-style states and goals expressed in the tradition of PDDL \citep{mcdermott1998pddl} and evaluated with an executable bounded-search solver in place of Fast Downward \citep{helmert2006fastdownward}. QuestBench \citep{li2025questbench} is the closest downstream benchmark because it asks whether LLMs can request the missing information needed for reasoning tasks.

Relative to this literature, our novelty claim is modest. We do not propose a new clarification architecture. Instead, we audit whether the \emph{quality} of synthetic clarification supervision matters under controlled token budgets. The executed artifacts support that question mainly for planning, not for formal reasoning broadly.

\section{Method}
\label{sec:method}
\subsection{Study Design}
All runs use the same base model, optimizer family, and approximate token budget. We compare:
\begin{enumerate}[leftmargin=*]
  \item \textbf{Answer-only}: fully specified logic and planning examples only.
  \item \textbf{Noisy clarification}: fully specified examples plus underspecified examples whose hidden fact fails the clean validator.
  \item \textbf{Validator-clean clarification}: fully specified examples plus underspecified examples that satisfy the clean validator.
\end{enumerate}

The implementation fixes Qwen2.5-3B-Instruct as the base model and fine-tunes with LoRA rank 16, $\alpha=32$, dropout 0.05, learning rate $2\times10^{-4}$, weight decay 0.01, cosine scheduling, batch size 8, gradient accumulation 2, maximum length 512, and one epoch. The target modules are \texttt{q\_proj}, \texttt{k\_proj}, \texttt{v\_proj}, \texttt{o\_proj}, \texttt{gate\_proj}, \texttt{up\_proj}, and \texttt{down\_proj}. Runs use bf16 on an NVIDIA RTX A6000 with 47.41\,GB VRAM.

\subsection{Token-Matched Training Sets}
The training selector in \texttt{exp/shared/train\_eval.py} matches conditions by token count rather than by raw example count. The final budgets are 25{,}539 tokens for answer-only, 25{,}506 for noisy, and 25{,}517 for clean. This yields 238 answer-only examples and 202 examples each for noisy and clean. The answer-only set contains 119 logic and 119 planning items; the clean and noisy sets contain 102 logic and 100 planning items each.

\subsection{Logic Construction}
The logic pipeline uses ProofWriter \citep{tafjord2021proofwriter}. Fully specified examples ask the model to output \texttt{ANSWER: true}, \texttt{false}, or \texttt{unknown}. For underspecified examples, the code inspects proof traces and selects a candidate hidden fact only when the canonical proof for a true example consists of a single supporting triple. The clean logic setting therefore ends up much narrower than the original proposal: it is effectively restricted to one-hop proof-backed deletions. Noisy logic examples remove a different non-question sentence from the same context.

This construction is mechanically precise but conservative. It also likely explains part of the logic-side weakness: the executed logic supervision is much smaller in semantic variety than the proposal originally envisioned.

\subsection{Planning Construction}
The planning pipeline generates 4- or 5-block Blocksworld states and goals, then checks solvability with an exact bounded breadth-first search up to horizon $H=8$. Let $I$ be the full initial state, $G$ the goal facts, and $h$ a hidden fact candidate. In the executed artifacts, candidate hidden facts are restricted to \texttt{clear(x)} and \texttt{handempty}. An example is accepted as \emph{clean} when:
\begin{enumerate}[leftmargin=*]
  \item $(I, G)$ is solvable within horizon 8,
  \item the visible state plus $h$ is solvable,
  \item a proxy complement state constructed from $\neg h$ is unsolvable, and
  \item no other candidate in the same instance satisfies the same solvable/unsolvable flip.
\end{enumerate}

Noisy planning examples hide another candidate fact from the same solvable instance that fails this clean rule. Planning prompts are rendered in three verbalization families: A for training, B for validation, and C for test-time paraphrase robustness.

\subsection{Audits and Known Deviations}
The pre-registered human audit was not completed. The artifact therefore substitutes a proxy audit over 50 examples per condition and domain, plus residual-noise checks over 40 clean examples. Proxy precision is 1.000 for clean and 0.000 for noisy in both logic and planning, but these values come from executable heuristics rather than an external human reviewer. More importantly, the residual-noise planning check finds a longer-horizon instability rate of 1.000, meaning every sampled clean planning label changed under the broader horizon test. This sharply limits how strongly the planning clean/noisy comparison can be interpreted.

\section{Experiments}
\label{sec:experiments}
\subsection{Setup}
The prepared data contain 120 clean logic training examples and 120 clean planning training examples, with 40-example validation and test sets for each synthetic evaluation split. QuestBench contributes 1{,}150 Logic-Q test examples and 7{,}500 Planning-Q test examples. For the main claim we follow the pre-registered seed 11; seeds 23 and 47 are treated as exploratory replications.

We report:
\begin{itemize}[leftmargin=*]
  \item \textbf{QuestBench Planning-Q question accuracy}: whether the generated clarification question matches a valid question on planning items.
  \item \textbf{Synthetic planning clean question accuracy}: question accuracy on held-out clean Blocksworld examples.
  \item \textbf{Planning paraphrase robustness}: question accuracy on planning verbalization families A, B, and C.
  \item \textbf{Retention}: fully specified planning and logic accuracy.
  \item \textbf{Cost}: train examples, runtime, and peak VRAM.
\end{itemize}

Because logic question accuracy is zero in the main run and the planning residual-noise audit is unfavorable, we treat the results as pilot evidence rather than confirmation of the full hypothesis.

\subsection{Main Results}
Table~\ref{tab:main} reports the pre-registered seed-11 comparison. The clean condition is the strongest planning system on both QuestBench Planning-Q and synthetic clean planning evaluation. It reaches 0.636 on QuestBench Planning-Q, a +0.022 absolute improvement over answer-only and a +0.620 improvement over noisy supervision. On the held-out synthetic planning-clean set, clean reaches 0.450 while both baselines remain at 0.000. Figure~\ref{fig:questbench} visualizes the same comparison.

The result is not a simple clarification-versus-answering story. Answer-only is already strong on QuestBench Planning-Q (0.614), while noisy clarification collapses to 0.016. This suggests that low-quality clarification supervision can be actively harmful even under matched token budgets.

\begin{table*}[t]
\caption{Seed-11 main results. Higher is better for all accuracy metrics. Quest Logic-Q and logic clean question accuracy remain 0.000 for all three conditions in the main run and are omitted here for readability.}
\label{tab:main}
\centering
\small
\begin{tabular}{lcccccccc}
\toprule
Method & Quest Planning-Q & Planning Clean & Family A & Family B & Family C & Planning Full & Logic Full & Runtime (s) \\
\midrule
Answer-only & 0.614 & 0.000 & 0.000 & 0.000 & 0.000 & \textbf{0.775} & \textbf{0.775} & 19.169 \\
Noisy clarif. & 0.016 & 0.000 & 0.075 & 0.000 & 0.000 & 0.500 & 0.575 & \textbf{19.262} \\
Clean clarif. & \textbf{0.636} & \textbf{0.450} & \textbf{0.550} & 0.000 & \textbf{0.450} & 0.650 & 0.650 & 18.893 \\
\bottomrule
\end{tabular}
\end{table*}

\begin{figure}[t]
\centering
\includegraphics[width=0.78\linewidth]{figures/questbench_bars.png}
\caption{QuestBench transfer comparison from the stored figure artifact. The planning-side advantage of clean supervision is visible, while logic remains essentially at zero in the executed run.}
\label{fig:questbench}
\end{figure}

Bootstrap intervals from the stored analysis artifact reinforce the planning result. For QuestBench Planning-Q, clean has a bootstrap mean of 0.636 with a 95\% interval of [0.606, 0.663], answer-only has 0.614 [0.583, 0.642], and noisy has 0.016 [0.009, 0.024]. The clean-minus-answer delta is 0.022 [0.014, 0.031], while clean-minus-noisy is 0.620 [0.590, 0.648].

\subsection{Retention and Robustness}
The improvement does not come for free. Answer-only retains the best fully specified planning accuracy at 0.775, while clean falls to 0.650 and noisy to 0.500. The pre-registered retention rule required clean to stay within two absolute points of answer-only; the realized planning drop is 12.5 points, so this criterion fails.

Figure~\ref{fig:robustness} shows that the clean system's planning behavior is not uniformly robust across paraphrases. It performs well on family A (0.550) and family C (0.450), but family B falls to 0.000. The noisy model is weak on every family, with only 0.075 on A and near-zero elsewhere. This means the clean advantage is real in some surface forms, but it is not yet stable across all prompt variants.

\begin{figure}[t]
\centering
\includegraphics[width=0.78\linewidth]{figures/planning_robustness.png}
\caption{Planning paraphrase robustness from the stored figure artifact. Clean supervision improves family A and family C, but not family B, so the transfer effect is not uniformly paraphrase-stable.}
\label{fig:robustness}
\end{figure}

\subsection{Exploratory Multi-Seed Summary}
Although the main study was pre-registered as one seed, the workspace contains runs for seeds 23 and 47. Table~\ref{tab:multiseed} summarizes those exploratory results. The broader pattern still favors clean over noisy on planning, but the spread is non-trivial, especially for noisy QuestBench Planning-Q accuracy (0.351 $\pm$ 0.311). Clean synthetic planning-clean accuracy is unusually stable at 0.450 $\pm$ 0.000 across all three seeds.

\begin{table}[t]
\caption{Exploratory three-seed summary from the aggregated result files. Values are mean $\pm$ standard deviation over seeds 11, 23, and 47.}
\label{tab:multiseed}
\centering
\small
\begin{tabular}{lcccc}
\toprule
Method & Quest Planning-Q & Planning Clean & Planning Full & Logic Full \\
\midrule
Answer-only & 0.568 $\pm$ 0.048 & 0.000 $\pm$ 0.000 & \textbf{0.733 $\pm$ 0.095} & 0.700 $\pm$ 0.130 \\
Noisy clarif. & 0.351 $\pm$ 0.311 & 0.008 $\pm$ 0.014 & 0.508 $\pm$ 0.014 & 0.700 $\pm$ 0.109 \\
Clean clarif. & \textbf{0.602 $\pm$ 0.044} & \textbf{0.450 $\pm$ 0.000} & 0.542 $\pm$ 0.095 & \textbf{0.717 $\pm$ 0.063} \\
\bottomrule
\end{tabular}
\end{table}

\subsection{Label Quality and Feasibility}
Table~\ref{tab:labelquality} reports the audit-related quantities available in the artifacts. The clean/noisy precision gap is dramatic under the executable proxy audit, but the lack of a human audit prevents strong causal claims about semantic cleanliness. The residual-noise check is mixed: logic shows 0.000 alternative-proof rate on 40 sampled items, whereas planning shows 1.000 instability under a longer horizon on 40 sampled items.

\begin{table}[t]
\caption{Available label-quality evidence from the audit artifacts. Wilson intervals are reported exactly as stored.}
\label{tab:labelquality}
\centering
\small
\begin{tabular}{lcc}
\toprule
Statistic & Logic & Planning \\
\midrule
Clean proxy precision & 1.000 [0.929, 1.000] & 1.000 [0.929, 1.000] \\
Noisy proxy precision & 0.000 [0.000, 0.071] & 0.000 [0.000, 0.071] \\
Clean audit sample size & 50 & 50 \\
Noisy audit sample size & 50 & 50 \\
Residual-noise sample size & 40 & 40 \\
Residual-noise rate & 0.000 & 1.000 \\
\bottomrule
\end{tabular}
\end{table}

The study is nevertheless cheap to run. Figure~\ref{fig:budget} shows that actual wall-clock time was far below the originally budgeted eight hours: data preparation took 0.0027 hours, the proxy audit 0.0030 hours, pilot training 0.0046 hours, and the full training stage 0.0465 hours in the stored analysis artifact. The main seed training runs take roughly 19 seconds each and use 21.65--22.81\,GB peak VRAM.

\begin{figure}[t]
\centering
\includegraphics[width=0.78\linewidth]{figures/budget_vs_actual.png}
\caption{Budgeted versus actual stage time from the stored figure artifact. The run is highly feasible on a single RTX A6000, but that low cost partly reflects the final study's reduced scope.}
\label{fig:budget}
\end{figure}

\section{Discussion and Limitations}
\label{sec:discussion}
The clean-versus-noisy comparison appears meaningful in planning. Clean supervision consistently beats noisy supervision and, in the main seed, slightly exceeds answer-only on QuestBench Planning-Q while dramatically outperforming both baselines on the synthetic clean planning set. This supports the narrow claim that \emph{some} clarification transfer depends on supervision quality rather than on the presence of clarification examples alone.

At the same time, the pilot falls short of the stronger hypothesis from the proposal.

First, logic transfer is not established. In the main seed, QuestBench Logic-Q question accuracy and synthetic logic clean question accuracy are both 0.000 for all three systems. Second, planning retention fails the pre-registered threshold: clean loses 12.5 absolute points of fully specified planning accuracy relative to answer-only. Third, the required human audit was not run. Fourth, the planning residual-noise analysis is severe: every sampled clean example changes status under the longer-horizon check. That finding suggests the current clean label may depend too strongly on the chosen horizon proxy.

There are also implementation deviations from the original proposal. The planning validator uses a bounded BFS solver rather than Fast Downward, planning hidden facts are restricted to \texttt{clear(x)} and \texttt{handempty}, and the clean logic selector is effectively limited to single-triple proofs. These choices make the paper reproducible, but they narrow the scientific claim substantially.

\section{Conclusion}
This paper revisited a simple question: does validator-clean clarification supervision transfer to formal reasoning better than matched noisy supervision? In the executed artifacts, the answer is ``partly, in planning.'' On the pre-registered seed, clean supervision improves QuestBench Planning-Q question accuracy to 0.636 and synthetic clean planning question accuracy to 0.450, while noisy supervision collapses to 0.016 and 0.000 respectively. However, logic transfer remains absent, fully specified planning retention degrades from 0.775 to 0.650, the human audit is missing, and the planning residual-noise rate under a longer horizon is 1.000.

The right interpretation is therefore not a confirmed formal-reasoning result, but a repaired pilot. Cleanliness appears to matter, yet the current artifacts do not justify a broad cross-domain claim. The next version should restore the human audit, strengthen the planning validator beyond the current horizon-sensitive proxy, and recover a more faithful logic setting before claiming that clarification supervision transfers to formal reasoning in general.

\bibliographystyle{plainnat}
\bibliography{refs}

\end{document}
